\section{REVISÃO DE LITERATURA}

Esta seção apresenta os principais trabalhos relacionados ao Aprendizado Federado aplicado à detecção de fraudes bancárias, organizada em quatro temas: a evolução do Aprendizado Federado, os estudos sobre modelos baseados em árvore em ambientes distribuídos, as pesquisas específicas sobre detecção de fraudes com esse tipo de abordagem, e uma análise comparativa que posiciona este trabalho em relação ao que já foi publicado.

\subsection{Evolução do Aprendizado Federado}

O Aprendizado Federado foi formalmente introduzido por McMahan et al. \cite{mcmahan2017fedavg} com o algoritmo FedAvg (\textit{Federated Averaging}), que estabeleceu as bases para o treinamento distribuído de modelos de aprendizado de máquina preservando a privacidade dos dados. Nesse trabalho seminal, os autores demonstraram que é possível treinar modelos de alta qualidade em dados distribuídos por meio da combinação das atualizações locais, sem centralizar os dados brutos. Yang et al. \cite{yang2019fl} depois sistematizaram o campo, distinguindo entre Aprendizado Federado Horizontal (participantes com os mesmos tipos de dados, mas sobre pessoas diferentes) e Vertical (participantes com dados sobre as mesmas pessoas, mas com tipos de informação diferentes) --- o presente trabalho se enquadra no primeiro tipo.

A questão de como escalar esse tipo de sistema para muitos participantes foi abordada por Bonawitz et al. \cite{bonawitz2019scale}, que propuseram protocolos para sistemas com milhões de dispositivos. Konečný et al. \cite{konecny2016communication} investigaram como reduzir o volume de dados que precisa ser transmitido a cada rodada, usando técnicas de compressão --- um problema importante quando os modelos são grandes ou a conexão de rede é limitada.

Levantamentos mais recentes da área documentam esse crescimento. Liu et al. \cite{liu2024survey} apresentaram uma visão atualizada com diferentes formas de classificar as abordagens de Aprendizado Federado, identificando direções em aberto como personalização e eficiência. Chai et al. \cite{chai2024survey} argumentam que sistemas federados devem ser avaliados não apenas pela qualidade do modelo produzido, mas também pela eficiência de comunicação e pela segurança. Shi \cite{shi2023survey} documentou a transição do processamento centralizado para o distribuído, destacando que o Aprendizado Federado surgiu em resposta às regulamentações de privacidade e às dificuldades práticas de transferir grandes volumes de dados entre instituições.

\subsection{Modelos Baseados em Árvore em Ambientes Federados}

A adaptação de algoritmos de boosting para o ambiente federado é uma área de pesquisa ativa e diretamente relevante para dados financeiros em formato tabular. Cheng et al. \cite{cheng2021secureboost} introduziram o SecureBoost, um sistema que permite treinar modelos XGBoost de forma federada sem perda de precisão em relação ao treinamento centralizado. O SecureBoost usa técnicas de combinação segura para evitar o compartilhamento direto dos gradientes de treinamento, mantendo acurácia próxima a 90\% em dados de fraude. Esse trabalho é fundamental para validar a viabilidade de modelos de árvore em cenários federados com requisitos rigorosos de privacidade.

Li et al. \cite{li2023fedtree} propuseram o FedTree, um sistema otimizado para treinamento federado de árvores de decisão que demonstra ganhos expressivos de eficiência. Os autores mostraram que implementações bem planejadas podem ser muito mais rápidas do que abordagens iniciais, aspecto importante para detecção de fraude em tempo real. Li, Wen e He \cite{li2020practical} complementam essa linha ao demonstrar que é possível manter a eficiência do LightGBM em ambientes distribuídos.

A questão da comunicação eficiente foi abordada por Le et al. \cite{le2021fedxgboost} com o FedXGBoost, que usa combinação segura para transmitir apenas informações agregadas dos histogramas de treino, sem expor os gradientes individuais. Liu et al. \cite{liu2020fedforest} exploraram uma abordagem alternativa usando encriptação para garantir privacidade em conjuntos de árvores de decisão, com custo computacional mais alto.

Uma vantagem importante dos modelos baseados em árvore em relação a redes neurais é que eles são naturalmente mais resistentes a ataques de reconstrução de dados. Zhu, Liu e Han \cite{zhu2019deepleakage} demonstraram que é possível recuperar dados de treinamento a partir das atualizações transmitidas em redes neurais. Modelos de árvore, ao transmitir informações sobre divisões e histogramas, tornam essa reconstrução muito mais difícil \cite{liu2024survey}.

\subsection{Detecção de Fraudes Bancárias e Aprendizado Federado}

A aplicação de Aprendizado Federado à detecção de fraudes financeiras tem recebido atenção crescente. Farrukh et al. \cite{farrukh2025blockchain} realizaram uma revisão comparando aprendizado de máquina centralizado com federado para detecção de fraudes, concluindo que o federado representa a evolução necessária para o setor bancário: embora métodos centralizados ofereçam alta precisão, eles criam silos de dados e pontos únicos de falha incompatíveis com as exigências regulatórias atuais.

Kennedy, Hilal e Momeni \cite{kennedy2025financial} investigaram especificamente o papel do Aprendizado Federado na segurança financeira, classificando a detecção de fraude em abertura de conta como uma tarefa de "alta exposição regulatória" onde o uso de aprendizado federado não é apenas uma vantagem técnica, mas uma necessidade de conformidade. O estudo documenta que a colaboração entre bancos via Aprendizado Federado pode reduzir a taxa de falsos positivos em até 15\% em comparação com modelos treinados isoladamente.

Claus, John e John \cite{claus2025collaborative} focaram em redes financeiras colaborativas, argumentando que fraudes que "saltam" de banco para banco --- onde um fraudador usa um banco para criar credenciais que depois usa em outro --- só podem ser detectadas por modelos treinados de forma colaborativa. Lynch, Abdelmoniem e Gill \cite{lynch2024fraud} compararam diretamente abordagens centralizadas e distribuídas para detecção de fraude, observando que a abordagem federada, embora tenha uma leve desvantagem em precisão absoluta, oferece maior resistência a ataques de reconstrução de dados e atende a requisitos regulatórios que a centralizada não consegue cumprir.

\subsection{Dados Heterogêneos e Equidade}

Um dos principais desafios em sistemas de Aprendizado Federado é quando os dados de cada participante têm distribuições muito diferentes entre si. Zhu et al. \cite{zhu2021noniid} documentam como essa heterogeneidade pode degradar significativamente o desempenho de modelos globais. Li et al. \cite{li2020fedprox} propuseram o FedProx como resposta a esse problema, adicionando um termo de regularização que limita o quanto o modelo de cada cliente pode se afastar do modelo global. Karimireddy et al. \cite{karimireddy2020scaffold} desenvolveram o SCAFFOLD, que usa variáveis de controle para corrigir o desvio acumulado causado por dados heterogêneos. Ambas as abordagens são alternativas ao FedAvg original e podem ser exploradas em trabalhos futuros.

A questão de equidade algorítmica em sistemas federados foi investigada por Huang et al. \cite{huang2022fairness}, que estudaram especificamente a relação entre equidade e desempenho no Aprendizado Federado Horizontal. Ao distribuir dados entre clientes, surge um compromisso entre manter a precisão global e garantir que o modelo não trate grupos específicos de forma desfavorável. Zhang, Kou e Wang \cite{zhang2020fairfl} propuseram o FairFL para reduzir viés demográfico em modelos federados preservando a privacidade. Hardt, Price e Srebro \cite{hardt2016equality} formalizaram conceitos de equidade algorítmica como igualdade de oportunidade (mesma taxa de detecção correta para grupos diferentes). No contexto deste trabalho, a métrica de equidade adotada segue o artigo do BAF \cite{jesus2022baf}: o FPR Ratio, que mede a razão entre as taxas de falso positivo de grupos demográficos diferentes e indica se o modelo penaliza desproporcionalmente certos grupos com rejeições incorretas.

\subsection{Privacidade e Segurança}

Preservar a privacidade em sistemas federados vai além de simplesmente não centralizar os dados. Bonawitz et al. \cite{bonawitz2017secagg} propuseram protocolos de Combinação Segura (\textit{Secure Aggregation}) que garantem que o servidor central não veja as atualizações individuais dos clientes, apenas o resultado combinado. Geyer, Klein e Nabi \cite{geyer2017differential} investigaram o uso de Privacidade Diferencial --- técnica que adiciona ruído estatístico controlado às atualizações para garantir privacidade formal --- em sistemas federados. Wei et al. \cite{wei2020privacy} formalizaram a relação entre privacidade e qualidade do modelo, mostrando que quanto maiores as garantias de privacidade, maior tende a ser a perda de desempenho. Este trabalho não implementa essas técnicas adicionais de privacidade, focando na avaliação de desempenho e equidade dos algoritmos de árvore.

\subsection{Comparação com os Trabalhos Relacionados}

A Tabela \ref{tab:trabalhos_relacionados} apresenta uma comparação dos principais trabalhos relacionados com o presente estudo.

\begin{table}[H]
\centering
\small
\begin{tabularx}{\textwidth}{|l|X|X|}
\hline
\textbf{Trabalho} & \textbf{Contribuição Principal} & \textbf{Limitação / Lacuna} \\
\hline
SecureBoost \cite{cheng2021secureboost} & Framework para XGBoost federado sem perda de precisão & Não avalia LightGBM e CatBoost \\
\hline
FedTree \cite{li2023fedtree} & Sistema otimizado para árvores federadas & Foco em eficiência, não em fraude \\
\hline
Kennedy et al. \cite{kennedy2025financial} & Classificação de abordagens FL em finanças & Análise teórica, sem experimentos comparativos \\
\hline
Farrukh et al. \cite{farrukh2025blockchain} & Revisão ML centralizado vs FL para fraude & Não compara algoritmos de boosting \\
\hline
BAF Dataset \cite{jesus2022baf} & Dataset de referência para fraude em abertura de conta & Avalia apenas LightGBM centralizado \\
\hline
Huang et al. \cite{huang2022fairness} & Equidade em Aprendizado Federado Horizontal & Não aplica a detecção de fraude \\
\hline
\textbf{Este trabalho} & \textbf{Comparação de XGBoost, LightGBM e CatBoost federados no BAF com análise de equidade} & --- \\
\hline
\end{tabularx}
\caption{Comparação dos trabalhos relacionados}
\label{tab:trabalhos_relacionados}
\end{table}

Este trabalho se diferencia da literatura em três aspectos. Primeiro, realiza a primeira comparação sistemática entre XGBoost, LightGBM e CatBoost em ambiente federado, preenchendo uma lacuna onde a maioria dos estudos foca em um único algoritmo. Segundo, utiliza o BAF --- o conjunto de dados de referência mais recente e realista para fraude em abertura de conta --- estendendo sua avaliação para o contexto federado com análise de equidade. Terceiro, avalia o comportamento de longo prazo das estratégias de combinação de modelos ao longo de 50 rodadas, revelando um fenômeno de degradação progressiva na estratégia Bagging que não foi documentado em experimentos mais curtos disponíveis na literatura.
